% 三版共用的标题编号、目录条目与书签设置，由各版布局在设置字号之后 \input。
% 目录标题及目录前后的分页由各版负责：竖版用 tocloft 的标题设置，双栏版本见 shared/twocolumn.tex。
% 每次排目录前须把 \BookTOCPreviousLevel 清零：竖版在布局中挂钩，双栏版本在重定义的 \tableofcontents 中清零。

% 保留原文的 0xPQ 编号：只关闭章节标题的自动编号。
\setcounter{secnumdepth}{-1}
% 目录展开到三级；PDF 书签包含全部标题层级并默认全部展开。
\setcounter{tocdepth}{3}
\hypersetup{bookmarksdepth=5,bookmarksopen=true}
\renewcommand{\contentsname}{目录}
% 方案 A：一级沿用苹方 Semibold；二三级沿用正文宋体与 Latin Modern Roman。
\newlength{\BookTOCUnit}
\settowidth{\BookTOCUnit}{\normalfont\fontsize{10}{15}\selectfont 中}
% 每条目录装入完整段落盒：续行保持本级缩进，页码固定在末行右端。
% 条目间使用可丢弃的竖直间距，页首不残留标题前留白。
% 根据前后层级禁止父标题与首个子条目之间分页，不捆绑整个章节。
\makeatletter
\newcount\BookTOCPreviousLevel
\BookTOCPreviousLevel=0
\newcommand{\BookTOCEntry}[4]{%
  \par
  \ifnum\BookTOCPreviousLevel=1 \penalty10000
  \else\ifnum\BookTOCPreviousLevel=2
    \ifnum#1=3 \penalty10000\else\penalty0\fi
  \else\penalty0\fi\fi
  \begingroup
  \@tempdima=0pt
  \ifnum\BookTOCPreviousLevel=1 \advance\@tempdima by5pt\fi
  \ifnum#1=1 \advance\@tempdima by12pt\fi
  \vskip\@tempdima
  \nointerlineskip
  \vbox{%
    \hsize=\linewidth
    \normalfont\fontsize{10}{15}\selectfont
    \ifnum#1=1 \sffamily\bfseries\fontsize{12}{15}\selectfont
    \else\ifnum#1=2 \fontsize{10.5}{15}\selectfont\fi\fi
    \leftskip=#2\relax \rightskip=30pt
    \parfillskip=-\rightskip \parindent=0pt \parskip=0pt
    \interlinepenalty=10000
    \noindent\strut#3\nobreak
    \ifnum#1=1 \hfill
    % 9pt 是实际点中心间距；两端各留 3pt，空间不足一格时不绘点。
    \else{\normalfont\fontsize{10}{15}\selectfont\color{BookTOCDots}%
      \kern3pt\leaders\hbox to9pt{\hss.\hss}\hfill\kern3pt}\fi
    \nobreak\hb@xt@20pt{\hfil\normalfont
      \ifnum#1=1 \fontsize{11}{15}\else\fontsize{10}{15}\fi\selectfont#4}\strut\par
  }%
  \endgroup
  \global\BookTOCPreviousLevel=#1\relax}
\renewcommand*{\l@section}[2]{%
  \ifnum\c@tocdepth>0 \BookTOCEntry{1}{0pt}{#1}{#2}\fi}
\renewcommand*{\l@subsection}[2]{%
  \ifnum\c@tocdepth>1 \BookTOCEntry{2}{2\BookTOCUnit}{#1}{#2}\fi}
\renewcommand*{\l@subsubsection}[2]{%
  \ifnum\c@tocdepth>2 \BookTOCEntry{3}{4\BookTOCUnit}{#1}{#2}\fi}
\makeatother
